\documentclass[usenames,dvipsnames]{beamer}

\mode<presentation> {
\usetheme{Montpellier}
}

\usepackage{graphicx} % Allows including images
\usepackage{booktabs} % Allows the use of \toprule, \midrule and \bottomrule in tables
\usepackage[frenchb]{babel}
\usepackage[T1]{fontenc}
\usepackage[utf8]{inputenc}
\setcounter{tocdepth}{2}
\graphicspath{{../Figures/}}
\DeclareMathOperator*{\argmin}{arg\,min}
\DeclareMathOperator*{\argmax}{arg\,max}
\newtheorem*{proposition}{Proposition}

% exercices comptés
\newcounter{numexos}%Création d'un compteur qui s'appelle numexos
\setcounter{numexos}{0}%initialisation du compteur
\newcommand{\exercice}[1]{%Création d'une macro ayant un paramètre
\addtocounter{numexos}{1}%chaque fois que cette macro est appelée, elle ajoute 1 au compteur numexos
\textcolor{DarkOrchid}{Exercice\,\thenumexos\,:}\,#1%Met en rouge Exercice et la valeur du compteur appelée par \thenumeexos
}

\setbeamertemplate{title page}{
    \centering

    \huge\bfseries
    Fondamentaux théoriques du machine learning

\begin{figure}[!h]
\includegraphics[width=8cm]{optimal_regression}
\end{figure}

    }

\title[FTML]{title}


\date{} % Date, can be changed to a custom date

\makeatletter
\let\@@magyar@captionfix\relax
\makeatother

\begin{document}

\begin{frame}
\titlepage % Print the title page as the first slide
\end{frame}

\begin{frame}
\frametitle{Overview of lecture 2} % Table of contents slide, comment this block out to remove it
\tableofcontents % Throughout your presentation, if you choose to use \section{} and \subsection{} commands, these will automatically be printed on this slide as an overview of your presentation
\end{frame}

\section{Mathematical toolbox for ML (part I)}%
\frame{\tableofcontents[currentsection]}
\label{sec:mathematical_tools_for_ml}

\begin{frame}{Mathematical toolbox}
    \begin{itemize}
        \item The aim of the course if to give an introduction to \textbf{{fundamental principles}} in ML.
	\item To do so, we will need an adapted mathematical toolbox and a bag of important results.
    \end{itemize}

    Why are mathematical aspects useful ? 
    \begin{itemize}
    	\item they allow a good comprehension of some theoretical results on ML
	\item these results allow a good choice of algorithms on practical problems (hopefully fast, accurate, etc.)
    \end{itemize}
    This section will give you an overview of the tools that will make you benefit more from the course if you are comfortable with them.
\end{frame}

\subsection{Linear algebra}%
\label{sub:linear_algebra}

\begin{frame}{Matricial calculus}
In machine learning, optimization or statistics we often write the inner product of two vectors of $ \mathbb{R}^d$ as a product of matrices. If $x\in \mathbb{R}^d$ writes:
\begin{equation*}
x=\begin{pmatrix}
x_{1}\\
...\\
x_{i}\\
...\\
x_{d}
\end{pmatrix}
\end{equation*}

And (with $T$ denoting the transposition),
\begin{equation*}
y^T=\begin{pmatrix}
y_{1}, \dots y_j, \dots, y_d\\
\end{pmatrix}
\end{equation*}

Then we have that

\begin{equation*}
    \langle x, y \rangle = y^Tx=x^Ty
\end{equation*}
\end{frame}

\subsection{Metrics}

\begin{frame}{Metrics}

    Let $D = \{x_1, \dots, x_n\}\subset \mathcal{X} $ be a dataset of $n$ samples, with labels $\{y_1, \dots, y_n\} \subset \mathcal{Y}$.

    There is a metric in the input space $ \mathcal{X}$ and in the output space $ \mathcal{Y}$.
    \begin{itemize}
    	\item The \textbf{metric} in $ \mathcal{X}$ determines to what extent two samples $x_i$ and $x_j$ should be considered similar or dissimilar.
	\item The \textbf{metric} in $ \mathcal{Y}$ determines to what extent two labels $y_i$ and $y_j$ should be considered similar or dissimilar.
    \end{itemize}

    This is very important during the complete processing of the data.
\end{frame}



\begin{frame}{Metrics in output space}

    A \textbf{{loss function}} $l$ is a map that measures the discrepancy between to elements of a set (for instance of a linear space).
$$
l: \left\{
    \begin{array}{ll}
	\mathcal{Y} \times \mathcal{Y} \rightarrow \mathbb{R}_+ & \\
	(y,z) \mapsto l(y,z)& 
    \end{array}
\right.
$$
    Typically, $z$ can represent our prediction for a given input $x$, $z= \tilde{f}(x)$, and $y$ the correct label.

\end{frame}

\begin{frame}{Most common loss functions}
%
	\begin{itemize}

			\item "0-1" loss for \textbf{binary classification}: $ \mathcal{Y}  = \{0, 1\}$ or $ \mathcal{Y}  = \{-1, 1\}$.

	\begin{equation}
	    l(y,z) = 1_{y\neq z}
	\end{equation}


		\item square loss for \textbf{{regression}}: $ \mathcal{Y}  = \mathbb{R} $. 

	\begin{equation}
	    l(y,z) = (y-z)^2
	\end{equation}


		\item absolute loss for \textbf{{regression}}: $ \mathcal{Y}  = \mathbb{R} $. 

	\begin{equation}
	    l(y,z) = |y-z|
	\end{equation}

	\end{itemize}
\end{frame}
\begin{frame}

    In \textbf{unsupervised learning}, there is notion of \textbf{output space !} (this might depend on the point of view)

\end{frame}

\begin{frame}
    \frametitle{Geometric distances in $\mathbb{R}^p $}
    $x=(x_1,...,x_p)$ and $y=(y_1,...,y_p)$ are $p$-dimensional \textbf{vectors}.
    \begin{itemize}
	\item $L_2$ : $||x-y||_2=\sqrt{\sum_{k=1}^{p} (x_k-y_k)^2}$ (Euclidian distance, $2$-norm distance)
	\item $L_1$ : $||x-y||_1=\sum_{k=1}^{p} |x_k-y_k|$ (Manhattan distance, $1$-norm distance)
	\item weighted $L_1$ :  $\sum_{k=1}^{p} w_k|x_k-y_k|$, wich each $w_k >0$
	\item $L_{\infty}$ : $\max ( x_1, \dots, x_n)$ (infinity norm distance, Chebyshev distance)
    \end{itemize}
\end{frame}


\begin{frame}{Choice of the metric}
    In some contexts, some usual metrics such as $L2$ might not be meaningful !

    \begin{figure}[htpb]
        \centering
        \includegraphics[width=0.8\linewidth]{supelec}
	\caption{In $ \mathbb{R}^{64}$, those three points form an equilateral triangle, \cite[]{Fix}}
        \label{fig:supelec}
    \end{figure}
\end{frame}

\begin{frame}
   \exercice{}Using \textbf{simulations/lecture\textunderscore 2/metrics/geometric\textunderscore data/build\textunderscore graphs.ipynb}, choose the metric and the threshold so that this graph (and the ones on the next slides) are built.

    \begin{figure}[htpb]
    	\centering
    	\includegraphics[width=0.6\textwidth]{thres_0_3_dist_weighted_manhattan}
    \end{figure}
\end{frame}

\begin{frame}{}
    \begin{figure}[htpb]
    	\centering
    	\includegraphics[width=0.8\textwidth]{thres_0_4_dist_weighted_manhattan}
    	\label{fig:thres_0_4_dist_weighted_manhattan}
    \end{figure}
\end{frame}

\begin{frame}{}
    \begin{figure}[htpb]
    	\centering
    	\includegraphics[width=0.8\textwidth]{thres_0_55_dist_infinie}
    	\label{fig:thres_0_55_dist_infinie}
    \end{figure}
\end{frame}

\begin{frame}{Non-geometric data}

    Not all data are geometric !

\end{frame}

\begin{frame}
    \frametitle{Hamming distance}
    \begin{itemize}
	\item  $\#\{x_i\neq y_i\}$ (Hamming distance)
	\item Levenshtein distance for strings (allows deletions and additions)
    \end{itemize}
\end{frame}

\begin{frame}
    \frametitle{General definition of a distance}
    A \textbf{distance}  on a set $E$ is an application $d : E\times E \rightarrow \mathbb{R}_+$ that must : 
    \begin{itemize}
	\item be \textbf{symetric} : $\forall x,y, d(x,y)=d(y,x)$
	\item \textbf{separate the values} : $\forall x,y, d(x,y)=0\Leftrightarrow x=y$
	\item	respect the \textbf{triangular inequality} $\forall x,y,z, d(x,y)\leq d(x,z)+d(y,z)$
    \end{itemize}
\end{frame}

\begin{frame}{Similarities}
    Sometimes, it is not possible to define a proper \textbf{distance} in the input space $ \mathcal{X}$ ! This may happen for instance is $ \mathcal{X}$ is a dataset of texts.
    \begin{itemize}
	\item When distances are unavailable, we can use \textbf{{Similarities}} or \textbf{{Dissimilarity}} to compare points.
	\item Dissimilarites are more general and don't always abide by the
	    distance axioms.
	\item Other examples: Adjacency in an oriented graph, custom agregated score to compare data. 
    \end{itemize}
\end{frame}

\begin{frame}{Example: cosine similarity}

    The \textbf{{cosine similarity}} may be used to compare texts.

    If $u$ and $v$ are vectors,
    \begin{equation}
	S_C(u, v) = \frac{(u|v)}{||u|| ||v||} 
    \end{equation}

    \begin{itemize}
	\item the \textbf{{bag of words representation}} allows us to build a vector from a text (one hot encoding).
	\item \textbf{{cosine\textunderscore similarity/scraper.py}} 
	\item \textbf{{cosine\textunderscore similarity/similarity.py}} 
    \end{itemize}
\end{frame}


\begin{frame}{Hybrid data}

    Sometimes each sample contains both numerical data and non-numerical data (text, categorical data.)

    See \textbf{hybrid\textunderscore data/}

    This is often the case in machine learning applications ! (database of customers, database of cars, etc.)

\end{frame}



\subsection{Statistics, probability theory}%
\label{sub:statistics_probability_theory}

\begin{frame}{Moments of a distribution}
\begin{definition}{Moments of a distribution}

    Let $X$ be a real random variabe on a probability space $\Omega$, and $k\in \mathbb{N}^*$. $X$ is said to have a moment of order $k$ if $E(|X|^k)<+\infty$, which means that:
    \begin{itemize}
	\item if $X$ is discrete, with image $X(\Omega)=(x_i)_{i\in \mathbb{N} }$, the series

	    \begin{equation*}
		\sum (x_i)^kP(X=x_i)
	    \end{equation*}

	    is \textbf{{absolutely}}  convergent. The moment is then equal to the sum of that series (without absolute value).
    \end{itemize}
\end{definition}
\end{frame}


\begin{frame}{Moments of a distribution}
\begin{definition}{Moments of a distribution}

    Let $X$ be a real random variabe on a probability space $\Omega$, and $k\in \mathbb{N}^*$. $X$ is said to have a moment of order $k$ if $E(|X|^k)<+\infty$, which means that:
    \begin{itemize}
	\item if $X$ is continuous with density $p(x)$, the integral
	    \begin{equation*}
	        \int^{+\infty}_{-\infty} x^kf(x)  dx 
	    \end{equation*}
	    is \textbf{{absolutely}}  convergent. The moment is then equal to the integral (without absolute value).
    \end{itemize}
\end{definition}
\end{frame}

\begin{frame}{Moments of a distribution}
\begin{proposition}

    Let $k_1<k_2$ be integers. Let $X$ be a real random variable. Then if $X$ has a moment of order $k_2$, $X$ also has a moment of order $k_1$.
    
\end{proposition}
\end{frame}


\begin{frame}{Moments of a distribution}
    \exercice{\textbf{Prove the proposition}}
\begin{proposition}

    Let $k_1<k_2$ be integers. Let $X$ be a real random variable. Then if $X$ has a moment of order $k_2$, $X$ also has a moment of order $k_1$.
    
\end{proposition}
\end{frame}


\begin{frame}{Expected value, variance}

\begin{definition}{Expected value, variance}

    \begin{itemize}
        \item If $X$ has a moment of order $1$, it is called the \textbf{{expected value}} 
	\item If $X$ has a moment of order $2$, then $X-E(X)$ also has a moment of order $2$. This moment is called the variance of $X$.
	    \begin{equation*}
		V(X)=E\big((X-E(X))^2\big)
	    \end{equation*}
	    We often note $\sigma(X)= \sqrt{Var(X)}$ (standard deviation).
    \end{itemize}
\end{definition}

\end{frame}

\begin{frame}{Expected value, variance}
\begin{proposition}
    \label{prop:var}

    Let $a$ and $b$ be real numbers, and $X$ a random variable that admits a moment of order $2$. Then
    \begin{equation*}
	Var(aX+b)=a^2Var(X)
    \end{equation*}
\end{proposition}
\end{frame}

\begin{frame}{Independence}
\begin{proposition}
    Let $(X_1, \dots, X_n)$ be $n$ mutually independent real random variables. Then if they all admit a moment of order $1$, then the product $X_1X_2\dots X_n$ also does admit a moment of order $1$ and 
    \begin{equation*}
	E(X_1X_2\dots X_n) = \prod^{n}_{i=1} E(X_i)
    \end{equation*}
    If they also admit moments of order $2$, then
    
    \begin{equation*}
	Var( \sum^{n}_{i=1} X_i ) = \sum^{n}_{i=1} Var(X_i)
    \end{equation*}
\end{proposition}
\end{frame}

\begin{frame}{Covariance}

    The covariance of the two real variables $X$ and $Y$ writes
    \begin{equation}
	Cov(X, Y) = E[(X-E[X])(Y-E[Y])]
    \end{equation}
\begin{lemma}
    Let $X, Y, Z\in \mathbb{R} $ be real random variables with a moment of order $2$. We have:
\begin{equation*}
    Cov(X+Y, Z) = Cov(X,Z)+Cov(Y,Z)
\end{equation*}

\begin{equation*}
    Var(X+Y) = Var(X)+Var(Y)+2Cov(X,Y)
\end{equation*}

\begin{equation*}
    |Cov(X,Y)|\leq \sigma(X)\sigma(Y)
\end{equation*}
\end{lemma}
    
\end{frame}

\begin{frame}{Convention}
    From now on, if we write $E(X)$ or $Var(X)$, we implicitely assume that the quantities are correctly defined. 
\end{frame}

\begin{frame}{Random vectors}
\begin{definition} 
    Let $X\in \mathbb{R}^d$ be a random vector.
\begin{equation*}
X=\begin{pmatrix}
X_{1}\\
...\\
X_{i}\\
...\\
X_{d}
\end{pmatrix}
\end{equation*}

The \textbf{{expected value}} of the vector writes
\begin{equation*}
    E(X)=\begin{pmatrix}
    E[X_{1}]\\
...\\
    E[X_{i}]\\
...\\
    E[X_{d}]
\end{pmatrix}
\end{equation*}

    The \textbf{{variance matrix}} (or \textbf{{covariance matrix, variance-covariance, dispersion matrix}})  $Var(X)$ is defined as

    \begin{equation*}
	[Var(X)]_{ij} = Cov(X_i, X_j)       
    \end{equation*}
\end{definition}
\end{frame}

\begin{frame}{Random vectors}
\begin{definition}

\begin{equation*}
X=\begin{pmatrix}
X_{1}\\
...\\
X_{i}\\
...\\
X_{d}
\end{pmatrix}
\end{equation*}

    The \textbf{{variance matrix}} (or \textbf{{covariance matrix, variance-covariance, dispersion matrix}})  $Var(X)$ is defined as

    \begin{equation*}
	[Var(X)]_{ij} = Cov(X_i, X_j)       
    \end{equation*}
\end{definition}
\end{frame}

\begin{frame}{Random vector}
    \exercice{\textbf{Random vector}}

    What does it mean to have a vector such that 

    \begin{equation}
	Var(X) = \lambda I_d
    \end{equation}
    
    ?
\end{frame}

\begin{frame}{Expected value as a minimization}

    \exercice{\textbf{Expected value as minimization.}}

    Show that $E(X)$ is the value that minimizes the function

    \begin{equation}
	f(t)=E\big( (X-t)^2 \big)
    \end{equation}
\end{frame}

\begin{frame}{Markov inequality}
\begin{proposition}{Markov inequality}

    Let $X$ be a real non-negative random variable (variable aléatoire réelle positive), such that $E(|X|)<+\infty$. Let $a>0$. Then

    \begin{equation*}
	P(X\geq a)\leq \frac{E(X)}{a} 
    \end{equation*}
    
\end{proposition}
\end{frame}

\begin{frame}{Chebychev inequality}
\begin{proposition}{Chebyshev inequality}
    Let $X$ ba a real random variable, such that $E(|X|^2)<+\infty$. Let $a>0$. Then

    \begin{equation*}
	P(|X-E[X]|>a)\leq \frac{Var(X)}{a^2} 
    \end{equation*}
\end{proposition}
\end{frame}

\begin{frame}{Weak law of large numbers}
\begin{theorem}{Weak law of large numbers}

    Let $(X_n)_{n\in \mathbb{N} }$ be a sequence of i.i.d. variables that have a moment of order $2$. We note $m$ their expected value. Then

    \begin{equation*}
	\forall \epsilon > 0, \lim_{n\rightarrow+\infty}P\big( |\frac{1}{n} \sum^{n}_{i=1} X_i-m|\geq \epsilon \big)= 0
    \end{equation*}

We say that we have \textbf{{convergence in probability}}.
\end{theorem}
\end{frame}

\begin{frame}{Standard deviation of the average}
    If $E(S_n)=m$, then

    
    \begin{equation}
	\sqrt{Var\Big( S_n-m \Big)} = \frac{\sigma}{ \sqrt{n} } 
    \end{equation}
\end{frame}


\begin{frame}[allowframebreaks]
\frametitle{References}
\bibliographystyle{apalike}
\bibliography{../../epita.bib} 
\end{frame}

\end{document} 
